% 与 combinatorics-mod.md 同内容（2026-11-20）。xelatex 编译；Overleaf 上把 Songti SC / Heiti SC 换成 FandolSong / FandolHei
\documentclass[10pt]{article}
\usepackage{xeCJK}
\setCJKmainfont{Songti SC}[BoldFont=Heiti SC]
\setCJKmonofont{Heiti SC}
\setmonofont{Menlo}[Scale=0.9]
\usepackage[a4paper,margin=1.6cm]{geometry}
\usepackage{amsmath,amssymb}
\usepackage{booktabs,array,tabularx}
\usepackage[table]{xcolor}
\usepackage{listings}
\usepackage{enumitem}
\usepackage{titlesec}
\usepackage{fancyvrb}

\definecolor{codefg}{RGB}{150,30,60}
\definecolor{codebg}{gray}{0.95}
\definecolor{note}{gray}{0.35}
\newcommand{\code}[1]{\texttt{\textcolor{codefg}{#1}}}
\lstset{language=Python,basicstyle=\ttfamily\small,backgroundcolor=\color{codebg},
  frame=none,xleftmargin=4pt,aboveskip=4pt,belowskip=4pt,columns=fullflexible,
  keywordstyle=\color{blue!60!black},commentstyle=\color{note},showstringspaces=false}
\titleformat{\section}{\large\bfseries}{\thesection.}{0.5em}{}
\titlespacing*{\section}{0pt}{10pt}{4pt}
\setlist{nosep,leftmargin=1.2em}
\renewcommand{\arraystretch}{1.3}
\setlength{\parindent}{0pt}
\setlength{\parskip}{3pt}
\pagestyle{plain}

\begin{document}

{\LARGE\bfseries 组合 / 取模 速查}\par
{\small\color{note} 2026-11-20 整理自手写两页 + 1866 / 1569 的讨论。展开版见 blog：advanced/modular.md · advanced/combinatorics.md}

\setcounter{section}{-1}
\section{总览：工具链 + 三种化法}

\begin{Verbatim}[fontsize=\small,frame=single,framesep=6pt]
工具链（前一个是后一个的前置）
  快速幂 (50) ──► 取模规则 (372) ──► 逆元 y^(p-2) ──► C(n,k) mod p (2400)

应用：计数题的三种化法
  ├─ 化成公式      2929  隔板 + 容斥                不取模
  ├─ 化成 DP       1866  第一类斯特林数              取模，不用 C
  └─ 结构 + 公式   1569  递归里每层乘 C(L+R, L)       C 取模
                   1916  多叉树拓扑序 n!/∏size        阶乘 + 逆元
\end{Verbatim}

共同点：问「有多少种」，答案大，模 $10^9+7$；不枚举方案，直接算个数。

\section{Fast pow}

$x^{13} = x^8 \cdot x^4 \cdot x^1$（$13 = 1101_2$）。$x^1, x^2, x^4, x^8$ 每个是前一个的平方 $\to$ 缓存着往上平方（via cached result），只乘指数二进制里是 1 的那几位。$O(\log n)$。

\begin{lstlisting}
def fast_pow(base, exp, mod):
    ans = 1
    base %= mod
    while exp > 0:
        if exp & 1:                    # 这一位是 1 才乘进去
            ans = ans * base % mod
        base = base * base % mod       # x^(2^i) -> x^(2^(i+1))
        exp >>= 1
    return ans
\end{lstlisting}
Python 直接 \code{pow(base, exp, mod)}。

\section{mod 规则：取模能放在哪}

\begin{tabularx}{\linewidth}{@{}l l X@{}}
\toprule
& \textbf{能先取模吗} & \textbf{怎么办}\\
\midrule
a) 底数 & $\checkmark$\ $(a \bmod m)^k \equiv a^k$ & \code{base \%= mod}\\
b) 乘积（加法同理） & $\checkmark$\ $xy \bmod m = (x \bmod m)(y \bmod m) \bmod m$ & 每步都模\\
c) 指数 & $\times$\ \ $a^b \not\equiv a^{b \bmod m}$ & 372：按十进制位递推；或满足条件时模 $\varphi(m)$\\
d) 除法 & $\times$\ \ $\frac{x}{y} \bmod m \ne \frac{x \bmod m}{y \bmod m}$ & 乘逆元：$\frac{x}{y} \bmod m = x \cdot y^{-1} \bmod m$\\
\bottomrule
\end{tabularx}

\begin{itemize}
\item 题目要的永远是\textbf{结果}取模；c) 是说「先把指数模小」这条捷径走不通。反例（$m=5$）：$2^5 \bmod 5 = 2$，而 $2^{5 \bmod 5} = 1$
\item d) 反例（$m=5$）：$12/4 = 3$；先模得 $2/4$，不是整数
\item \textbf{372}（$a^b \bmod 1337$，$b$ 是 base 10 的数组）：读数字是 \code{x = 10x + d}，搬到指数上 $a^{10x+d} = (a^x)^{10} \cdot a^d$，每读一位 \code{ans = ans\^{}10 · a\^{}d}，马上模
\item 十进制 = decimal / base 10；快速幂 = binary exponentiation
\end{itemize}

\section{逆元 inverse $y^{-1}$}

\textbf{定义}：$y \cdot y^{-1} \equiv 1 \pmod m$，模下的「$1/y$」。

\textbf{求法}：$m$ 是质数、$y \bmod m \ne 0$ 时，$y^{-1} \equiv y^{m-2} \pmod m$

\textbf{证明}：
\begin{enumerate}
\item 费马小定理：$m$ 是质数、$y \bmod m \ne 0 \Rightarrow y^{m-1} \equiv 1$
\item $\Rightarrow y \cdot y^{m-2} \equiv 1 \Rightarrow$ 按定义 $y^{-1} = y^{m-2}$
\end{enumerate}

\textbf{费马小定理怎么来（剩余系重排法）}：$y \cdot \{1, \dots, m-1\}$ 模 $m$ 后只是 $\{1, \dots, m-1\}$ 换了顺序。两边相乘 $y^{m-1}(m-1)! \equiv (m-1)!$，约掉 $(m-1)!$。

\begin{center}
\begin{tabular}{@{}l cccc@{}}
\toprule
$m=5,\ y=2$ & 1 & 2 & 3 & 4\\
\midrule
$2i \bmod 5$ & 2 & 4 & 1 & 3\\
\bottomrule
\end{tabular}
\end{center}

\begin{itemize}
\item 「只是重排」要两条：不出现 0、不撞车，都靠「质数整除乘积 $\Rightarrow$ 整除某个因子」
\item 非质数就坏：$m=4,\ y=2 \to 2, 0, 2$。372 的 $1337 = 7 \times 191$ 不是质数，用不了
\item 其他证法：二项式 + 归纳（$C_p^i$ 是 $p$ 的倍数）、项链计数、拉格朗日定理
\end{itemize}

\textbf{代码}：\code{pow(y, MOD - 2, MOD)}，或 3.8+ 的 \code{pow(y, -1, MOD)}。

\section{Combination via mod}

\[
C_n^k = \frac{n!}{k!\,(n-k)!}
\]
select $k$ from $n$ regardless of order：先按顺序挑 $\frac{n!}{(n-k)!}$ 种，同一组被数了 $k!$ 遍，除掉。$C_n^0 = C_n^n = 1$，$C_n^k = C_n^{n-k}$，$k > n$ 为 0。

\textbf{算法：$C_n^k \bmod m$}。直接套公式要么数太大，要么有除法不能边算边模 $\to$ 三种办法：

\begin{tabularx}{\linewidth}{@{}l X X X@{}}
\toprule
& \textbf{杨辉三角}（comb table） & \textbf{阶乘 + 逆元}（fact + inv table） & \textbf{API} \code{math.comb}\\
\midrule
怎么绕开除法 & 只用加法 & 除法换乘逆元 & 大整数精确算，最后 \code{\% MOD}\\
预处理 / 查询 & $O(n^2)$ / $O(1)$ & $O(n)$ / $O(1)$ & 无 / 每次 $O(k)$ 大整数\\
适用 & $n \le$ 几千；模数不是质数也行 & $n$ 到 $10^5 \sim 10^6$；模数必须是质数 & 只算几次、$n$ 几千内；只有 Python\\
\bottomrule
\end{tabularx}

\textbf{杨辉三角}：盯第 $a$ 个元素，选它 / 不选它
\[
C_a^b = C_{a-1}^{b-1} + C_{a-1}^{b},\qquad C_a^0 = 1
\]
\begin{lstlisting}
comb = [[0] * (n + 1) for _ in range(n + 1)]
for a in range(n + 1):
    comb[a][0] = 1
    for b in range(1, a + 1):
        comb[a][b] = (comb[a-1][b-1] + comb[a-1][b]) % MOD   # 括号别漏
\end{lstlisting}

\textbf{阶乘 + 逆阶乘}：只求一次逆元，$(i-1)!^{-1} = i!^{-1} \cdot i$ 从后往前推
\begin{lstlisting}
fact = [1] * N
for i in range(1, N):
    fact[i] = fact[i - 1] * i % MOD
inv_fact = [1] * N
inv_fact[N - 1] = pow(fact[N - 1], MOD - 2, MOD)
for i in range(N - 1, 0, -1):
    inv_fact[i - 1] = inv_fact[i] * i % MOD

def C(n, r):
    if r < 0 or r > n:
        return 0
    return fact[n] * inv_fact[r] % MOD * inv_fact[n - r] % MOD
\end{lstlisting}

\section{Usage}

\begin{tabularx}{\linewidth}{@{}l X X@{}}
\toprule
\textbf{题} & \textbf{化成什么} & \textbf{用到}\\
\midrule
50 Pow(x, n) & — & 快速幂模板\\
372 Super Pow & 按十进制位递推 & 快速幂 + 规则 c)\\
2400 恰好 $k$ 步 & $r = \frac{k+d}{2}$，答案 $C_k^r$；$|d| > k$ 或 $k+d$ 奇数为 0 & C 取模（当模板练：把 \code{math.comb} 换成阶乘表）\\
2929 分糖有上限 & 隔板 $C_{n+2}^2$ + 容斥 & 不取模，$k=2$ 直接写\\
1866 看得见 $k$ 根 & 盯最矮那根：\code{dp[i][j] = dp[i-1][j-1] + (i-1)·dp[i-1][j]} & 取模，杨辉三角形状的 DP，不用 C\\
1569 同一棵 BST & \code{ways = C(L+R, L) · ways(左) · ways(右)}，答案 $-1$ & C 取模（杨辉三角）\\
1916 蚁群建房 & $\frac{n!}{\prod \mathrm{size}(v)}$ & 阶乘 + 逆元\\
\bottomrule
\end{tabularx}

\textbf{纯模板练习}：3405（$C_{n-1}^k \cdot m \cdot (m-1)^{n-1-k}$，最推荐）· 62 / 1641（热身，不取模）· 1621 · 1735（多查询）· 2954（多重组合数）

\section{坑}
\begin{itemize}
\item \code{a + b * c \% MOD} 只模了乘积项：\code{\%} 和 \code{*} 同级、比 \code{+} 高。1866、1569 都踩了 $\to$ 整个式子加括号
\item \code{(ways - 1) \% MOD}：ways 模完可能是 0，直接减得负数
\item 1569 有序输入退化成链，递归 $\sim$1000 层，撞 Python 默认上限 $\to$ \code{sys.setrecursionlimit}
\item 逆元要求模数是质数；$10^9+7$ 是，1337 不是
\end{itemize}

\vspace{4pt}
{\small\color{note}
\textbf{手写稿里改了的地方}
\begin{itemize}
\item Fast pow 框里 \code{while i < n: ans *= x\^{}i; i <<= 1} 会把 $x^1 x^2 x^4 x^8$ 全乘进去，要只乘 $n$ 的二进制里是 1 的位（\code{if exp \& 1}）
\item reverse $\to$ \textbf{inverse}（逆元）；$y^{-1} \equiv y^{m-2}$ 的前提「$m$ 是质数」要写在公式旁边
\item b) 右边漏了最外层 mod $m$
\end{itemize}}

\end{document}
